\documentclass[oneside,14pt]{book}

\usepackage[T1,T2A]{fontenc}
\usepackage[utf8]{inputenc}
\usepackage[english,russian]{babel}
\usepackage{indentfirst}

%\usepackage[paperwidth=15cm,paperheight=7.5cm]{geometry} % планшет
%\usepackage[paperwidth=297mm,paperheight=210mm]{geometry} % A4
%\usepackage[paperwidth=148mm,paperheight=105mm]{geometry} % A5
\usepackage[paperwidth=18cm,paperheight=13cm,margin=5mm]{geometry} % экран*2

%\usepackage[colorlinks=true,
%]{hyperref}
%\newcommand{\email}[2]{#1\ \href{mailto:#2}{<\nolinkurl{#2}>}}
% %\newcommand{\email}[2]{\emph{#1\ <#2>}}
\usepackage[unicode,colorlinks,
pdftitle={Azbuka ARMaturschika (ru)},
pdfauthor={(c) Dmitry Ponyatov <dponyatov@gmail.com>, SSAU ASCL},
pdfsubject={ru manual on writing programs for Cortex-M MCUs},
pdfkeywords={ARM} {Cortex} {MCU} {ARMatura} {Arduino} {SSAU}
]{hyperref}

\usepackage{wrapfig}
\usepackage{graphicx}
\usepackage{epstopdf}
\DeclareGraphicsExtensions{.eps}

\usepackage{listings}
\usepackage{dirtree}
\usepackage[usenames,dvipsnames,svgnames]{xcolor}
\newcommand{\cppcolor}{\color[rgb]{0.94, 0.97, 1.0}} % Alice blue
\newcommand{\asmcolor}{\color[rgb]{0.98, 0.92, 0.84}} % Antique white
\newcommand{\ldcolor}{\color[rgb]{0.67, 0.9, 0.93}} % Blizzard Blue
\newcommand{\concolor}{\color[rgb]{0.88, 1.0, 1.0}} % Light cyan
\newcommand{\makecolor}{\color[rgb]{0.9, 0.9, 0.98}} % Lavender mist

%\definecolor{cppcolor}{rgb}{0.94, 0.97, 1.0}
\lstset{frame=single,
numbers=left, numberstyle=\small, numbersep=1mm,
tabsize=4,
keywordstyle=\color{blue}\texttt,
commentstyle=\color{cyan}\texttt,
inputencoding=utf8, extendedchars=true, showspaces=false  
%keepspaces=true,
}
\usepackage{lstlangarm}
\usepackage{lstlanggnumake}
\usepackage{lstlanggnuld}

\lstdefinestyle{cpp}{language=C++,backgroundcolor=\cppcolor}
\lstdefinestyle{asm}{language={[ARM]Assembler},backgroundcolor=\asmcolor}
\lstdefinestyle{gnuld}{language={[GNU]LinkScript},backgroundcolor=\ldcolor}
\lstdefinestyle{con}{backgroundcolor=\concolor}
\lstdefinestyle{mk}{language=[GNU]Make,backgroundcolor=\makecolor,showtabs=true}

\newcommand{\cm}[1]{Cortex-M#1}
\newcommand{\cx}{\cm{x}}

\newcommand{\vld}{STM32VLDISCOVERY}

%\renewcommand{\url}[1]{\textbf{#1}}
\newcommand{\email}[1]{$<$\href{mailto:#1}{\textbf{#1}}$>$}

\newcommand{\cpp}{$C^{+^{+}}$}

\newcommand{\cp}[1]{\footnote{копипаста: #1}}

\newcommand{\thmod}{Thumb}
\newcommand{\armod}{ARM}
\newcommand{\arm}{ARM}

\newcommand{\Reg}[1]{\textbf{#1}}
\newcommand{\R}[1]{\Reg{R#1}}

\newcommand{\periph}[1]{\texttt{#1}}
\newcommand{\jtag}{\periph{JTAG}}

\usepackage{wasysym} % smileys
\usepackage{gensymb} % celsius
\usepackage{amssymb} % windows key
\usepackage{textcomp} % bigcircle

\usepackage[os=win]{menukeys}
\newcommand{\winstart}{$\boxplus$}
\newcommand{\file}[1]{\textbf{\textsf{#1}}}
\newcommand{\window}[1]{\textbf{\textit{#1}}}
\newcommand{\alarm}[1]{{\color{DarkRed}#1}}
\newcommand{\wcmd}[1]{\keys{\winstart+R}\ \directory{#1}}
\newcommand{\checkbox}{$\boxtimes$}
\newcommand{\uncheckbox}{$\square$}
\newcommand{\lms}{\keys{$\lhd$}}
\newcommand{\rms}{\keys{$\rhd$}}
\newcommand{\eclpx}{\window{Project Explorer}}

\newcommand{\win}[1]{\includegraphics[height=10ex]{fig/winlogo.jpg} #1}
\newcommand{\lin}[1]{\includegraphics[height=10ex]{fig/linuxcolor.png} #1}
\newcommand{\bug}{\includegraphics[height=10ex]{fig/iconbug.png}}

\newcommand{\linux}{Linux}
\newcommand{\git}{Git}
\newcommand{\eclipse}{\textcircled{$\equiv$}\textsc{eclipse}}
%\newcommand{\term}[1]{\underline{#1}}
\newcommand{\term}[1]{\underline{\color{DarkBlue} #1}}
\newcommand{\miktex}{MiK\TeX}
\newcommand{\internet}{Internet}
\newcommand{\ql}{Quantum$^{\circledR}L^{e}aPs$}
\newcommand{\gdb}{GDB}

\newcommand{\make}{\file{make}}
\newcommand{\makefile}{\file{Makefile}}

\usepackage{tocloft}
\newcommand{\listlabname}{Лабораторные работы}
\newlistof{lab}{ex}{\listlabname}
\newcommand{\labpart}[1]{\addcontentsline{ex}{part}{#1}}
\newcounter{labworkcounter}
\newcommand{\labwork}[1]{
\refstepcounter{labworkcounter}
\section*{ЛР\thelabworkcounter: #1}
\addcontentsline{toc}{subsection}{ЛР\thelabworkcounter: #1}
\addcontentsline{ex}{section}{ЛР\thelabworkcounter: #1}
}
\newcommand{\labref}[1]{ЛР\ref{#1}}

\newcommand{\thetitle}{Азбука халтурщика-ARMатурщика}

\newcommand{\mytitle}[1]{
\title{\Huge{\thetitle}\\
#1\\
\normalsize{учебный курс по микроконтроллерам \cx:\\
Миландр 1986ВЕ, STM32F, LPC21xx}}
}

\author{(copypasta) Понятов Д.А. \email{dponyatov@gmail.com}, ИКП СГАУ}
